% Problem record op_6690dfacb75e8dc0. This TeX file is derived from
% database/problems_json/op_6690dfacb75e8dc0.json by scripts/sync-tex.mjs; edit the JSON record
% and rerun the script rather than editing this file.
\section{Quantum capacity of higher-dimensional depolarizing channels}
\label{sec:op_6690dfacb75e8dc0}

\paragraph{Problem.}
What is the unassisted quantum capacity of a depolarizing channel in every finite dimension $d\geq3$? For $0\leq\lambda\leq1$, define the channel on $\mathcal L(\mathbb C^d)$ by
\begin{equation}
 \mathcal N_{d,\lambda}(\rho)=\lambda\rho+(1-\lambda)\operatorname{Tr}(\rho)\frac{I_d}{d}.
 \label{eq:h14-channel}
\end{equation}
The capacity $Q(\mathcal N)$ is the supremum of asymptotic qubit rates achievable with entanglement fidelity tending to one by arbitrary encoders and decoders over independent uses, without entanglement or classical-communication assistance. For a complementary channel $\mathcal N^c$ from any Stinespring dilation and $S(\sigma)=-\operatorname{Tr}\sigma\log_2\sigma$, it equals
\begin{equation}
 Q(\mathcal N)=\sup_{n\geq1}\frac1n\max_{\rho\in\mathcal D((\mathbb C^d)^{\otimes n})}\left[S(\mathcal N^{\otimes n}(\rho))-S((\mathcal N^c)^{\otimes n}(\rho))\right].
 \label{eq:h14-capacity}
\end{equation}
Determine Eq.~\eqref{eq:h14-capacity} for the family in Eq.~\eqref{eq:h14-channel}, including the exact boundary between zero and positive capacity. The unresolved parameter regime is
\begin{equation}
 \frac{d+2}{2(d+1)}<\lambda<1.
 \label{eq:h14-range}
\end{equation}
An evaluation of one input state or one finite block length alone does not determine the supremum in Eq.~\eqref{eq:h14-capacity} throughout Eq.~\eqref{eq:h14-range}.

\subsection*{Status}

Unsolved

\subsection*{Source}

The supplied note derives this question from Hayashi's treatment of quantum capacity. Etxezarreta Martinez, deMarti iOlius, and Crespo explicitly state the unresolved capacity problem for the finite-dimensional depolarizing family in Section~I, p.~2, and define the channel in Section~II, Eq.~(3) \sourcecite{ref:h14-emc}{EMC23}. Their replacement-noise parameter is $p=1-\lambda$.

\subsection*{Progress}

\begin{itemize}
  \item The channel has zero capacity for $1-\lambda\geq d/[2(d+1)]$ because it is antidegradable; at $\lambda=1$ it has capacity $\log_2 d$. Put $p=1-\lambda$ and $q=p(1-d^{-2})$. The maximally mixed input gives the hashing lower bound, while the no-cloning bound gives
\begin{equation}
 \max\{0,\log_2 d-h_2(q)-q\log_2(d^2-1)\}\leq Q(\mathcal N_{d,\lambda})\leq\max\!\left\{0,1-\frac{2(d+1)p}{d}\right\}\log_2 d,
 \label{eq:h14-bounds}
\end{equation}
where $h_2(q)=-q\log_2q-(1-q)\log_2(1-q)$. These bounds in Eq.~\eqref{eq:h14-bounds} leave a gap in Eq.~\eqref{eq:h14-range}; see Section~II, Eqs.~(4)--(6) \sourcecite{ref:h14-emc}{EMC23}.

  \item Kianvash, Fanizza, and Giovannetti construct degradable flagged extensions giving improved upper bounds for qudit depolarizing channels. Their Section~5.4 and Figure~2 include a nontrivial dimension-four comparison; optimizing these extension bounds does not supply matching transmission codes \sourcecite{ref:h14-kfg}{KFG22}.

  \item In the limit of increasing dimension, the no-cloning and hashing bounds meet after division by $\log_2 d$: for fixed $p$, the normalized capacity tends to $\max\{0,1-2p\}$. Theorem~1 and Corollary~1 analyze this normalized gap. This asymptotic statement neither evaluates the capacity for a fixed finite $d\geq3$ nor proves that the absolute gap in qubits vanishes \sourcecite{ref:h14-emc}{EMC23}.
\end{itemize}

\subsection*{References}

\begin{enumerate}[label={},leftmargin=4.8em,labelsep=0.5em]
  \footnotesize
  \item[\textup{[EMC23]}]\label{ref:h14-emc}
  J. Etxezarreta Martinez, A. deMarti iOlius, and P. M. Crespo, ``Superadditivity effects of quantum capacity decrease with the dimension for qudit depolarizing channels,'' \emph{Physical Review A} \textbf{108}, 032602 (2023). \href{https://doi.org/10.1103/PhysRevA.108.032602}{doi:10.1103/PhysRevA.108.032602}; \href{https://arxiv.org/abs/2301.10132}{arXiv:2301.10132}.

  \item[\textup{[KFG22]}]\label{ref:h14-kfg}
  F. Kianvash, M. Fanizza, and V. Giovannetti, ``Bounding the quantum capacity with flagged extensions,'' \emph{Quantum} \textbf{6}, 647 (2022). \href{https://doi.org/10.22331/q-2022-02-09-647}{doi:10.22331/q-2022-02-09-647}; \href{https://arxiv.org/abs/2008.02461}{arXiv:2008.02461}.
\end{enumerate}

\subsection*{Comment}

Audited on 2026-09-09 against the full cited papers and subsequent capacity literature. The existing qubit Pauli-capacity record already contains the dimension-two depolarizing problem and its July--August~2026 coding progress; the present record retains only the higher-dimensional remainder. No existing general qudit-Pauli-capacity statement was found. The depolarizing channel's additive classical Holevo capacity is a different quantity and does not remove the coherent-information regularization in Eq.~\eqref{eq:h14-capacity}. The large-dimension result is in qudits per channel use, obtained by normalizing rates by $\log_2 d$, and supplies no fixed-dimension solution.

\subsection*{Field}

Quantum Communication

\subsection*{Topic}

Quantum capacity; Additivity and regularization

\subsection*{ID}

\texttt{op\_6690dfacb75e8dc0}
